\setchapter{1}{Maximum Likelihood Estimation}
\ChapterHeading

\section{Likelihood function and the ML principle}

Let $f(y;\theta_0)$ denote the probability density function of the random variable $y$ given $\theta_0$.
Our objective is to estimate the true parameter vector $\theta_0$.
The joint density of $n$ iid observations of y is
\begin{equation}
    f(y_1,\ldots,y_n|\theta_0) = \prod_{i=1}^{n}f(y_i;\theta_0)
\end{equation}
Now taking $f(y_i;\theta)$ as a function of $\theta$ given $y_1,\ldots,y_n$ and write
\begin{equation}
    L(\theta|y_1,\ldots,y_n) = \prod_{i=1}^{n}f(y_i;\theta)
\end{equation}
gives the \Term{likelihood function}, the likelihood that the population parameter is $\theta$ given the observed sample.

\subsection{ML principle}
The \Term{Maximum Likelihood (ML) principle} suggests that estimators of the unknwon parameters 
are obtained by \Emph{maximizing $L(\theta)$ with respect to $\theta$}. 
If the random variable $y$ is discrete, $\hat{\theta}$ yields the value of $\theta$ that maximizes
the probability of observing the sample that actually occurred.\\

Usually, $\hat{\theta}$ can be obtained by solving the likelihood equation
\begin{equation}
    \frac{\partial \log L(\theta)}{\partial \theta}|_{\hat{\theta}}=0
\end{equation}
where $\log L(\theta|y_1,\ldots,y_n) = \sum_{i=1}^{n}\log f(y_i;\theta)$.
Occationally the ML estimator is not unique and $\log L(\theta)$ may have only one global maximum,
but multiple local maxima.\\

The \Term{regularity conditions} are the following:
\begin{enumerate}
    \item The first three derivatives of $\log f(y|\theta)$ with respect to $\theta$ are continuous and finite
    for almost all $y$ and for all $\theta$.
    \item For all values of $\theta$, $|\frac{\partial^3\log f(y|\theta)}{\partial\theta_j\partial\theta_k\partial\theta_l}|$
    is limited by a function that has finite expectations.
    \item The domain of $y$ does not depend on $\theta$.
    \Info{For example of the uniform distribution, since the density function depends on $\theta_0$,
    the theorem does not hold.}
    \item $\theta$ is an interior point to the compact parameter space $\Theta$.
\end{enumerate}

\subsection{Bartlett Identities}
Define the \Term{score vector} $S(\theta)$ and the \Term{Hessian matrix} $H(\theta)$ as 
\begin{equation}
    S(\theta) = \frac{\partial\log L(\theta)}{\partial\theta} = \sum_{i=1}^{n}\frac{\partial\log f(y_i|\theta)}{\partial\theta}
    = \sum_{i=1}^{n}s_i(\theta)
\end{equation}
\begin{equation}
    H(\theta) = \frac{\partial^2\log L(\theta)}{\partial\theta\partial\theta'} 
    = \sum_{i=1}^{n}\frac{\partial^2\log f(y_i|\theta)}{\partial\theta\partial\theta'}
    = \sum_{i=1}^{n}H_i(\theta)
\end{equation}
\begin{itemize}
    \item The \Term{First Bartlett identity} states that $E[s_i(\theta_0)]=0$. Hence, $E[S(\theta_0)]=0$.
    \item The \Term{Second Bartlett identity} states that $Var[s_i(\theta_0)]=-E[H_i(\theta_0)]$.
    \begin{proof}
    \begin{equation*}
        Var\left[\frac{\partial\log f(y_i|\theta)}{\partial\theta}\right] = -E\left[\frac{\partial^2\log f(y_i|\theta_0)}{\partial\theta\partial\theta'}\right]
    \end{equation*}
    where $y_i$ is a random variable.
    \end{proof}
\end{itemize}
$Var[s_i(\theta)] = E[s_i(\theta_0)s_i(\theta_0)']$ defines the \Term{Fisher's information matrix}
denoted as $\mcI(\theta_0)$. Hence, the result $Var[s_i(\theta)] = E[s_i(\theta_0)s_i(\theta_0)']=-E[H_i(\theta_0)]$
is also called the \Term{information matrix identity}.
\begin{proof}
    \begin{align*}
        Var[s_i(\theta)] & = E[[s_i(\theta_0)-\underbrace{E[s_i(\theta_0)]}_{=0,\text{ First Bartlett}}]\times[s_i(\theta_0)-\underbrace{E[s_i(\theta_0)]}_{=0,\text{ First Bartlett}}]']\\
        & = E[s_i(\theta_0)s_i(\theta)']
    \end{align*}
\end{proof}

\section{Properties of MLE}
Under the assumed regularity conditions of MLE possesses the following properties:
\begin{enumerate}
    \item \Emph{Consistency} meaning $plim\hat{\theta} = \theta_0$
    \item \Emph{Asymptotic normality} meaning $\sqrt{n}(\hat{\theta}-\theta_0)\xrightarrow{d}\mcN(0,\mcI(\theta_0)^{-1})$.
    \Info{The goal is to estimate the $[\mcI(\theta_0)]^{-1}$ matrix (Fisher's information matrix)
    $\to$ need to go to a known distribution (normal).}
    \item \Emph{Asymptotic efficiency} meaning that if $\tilde{\theta}$ is a regular consistent asymptotically normal 
    estimator such that $\sqrt{n}(\tilde{\theta}-\theta_0)\xrightarrow{d}\mcN(0,\Omega)$,
    then $\Omega - [\mcI(\theta_0)]^{-1}$ is positive semi-definite, i.e.,
    under these RC, the MLE asymptotically achieves the \Term{Cramer-Rao} lower bound which is given by 
    $[\mcI(\theta_0)]^{-1}$.
    \Info{Information matrix is smaller than $\Omega$ in the matrix sense, and the information matrix is the lower bound.}
    \item \Emph{Invariance} meaning that if $c(\theta)$ is a continuous and continuously differentiable one-to-one function,
    the MLE of $\gamma = c(\theta_0)$ is $c(\hat{\theta})$.
\end{enumerate}

Now, don't forget that $\theta_0$ is not known, so we want to estiamte the information matrix.

\subsection{Estimators of the information matrix}
There are three commonly used estimators of $\mcI(\theta_0)$:
\begin{itemize}
    \item \Emph{Expected Information}: If the form of the expected values of the second derivatives 
    of the log-likelihood function is known, then we can evaluate $\mcI(\theta)$ at $\hat{\theta}$. 
    \Info{This rarely happens in practice.}
    \item \Emph{Observed information}: Simply use $-\hat{H}(\hat{\theta})$.
    \begin{proof}
    \begin{align*}
        \mcI(\theta_0) & = Var[s_i(\theta_0)] = -E[H(\theta_0)]\\
        & = Var\left[\frac{\partial\log f(y_i|\theta_0)}{\partial\theta_0}\right]\\
        & = -E\left[\frac{\partial^2\log f(y_i|\theta_0)}{\partial\theta\partial\theta'}\right]\\
        & \underset{\text{avg.}}{=} \frac{1}{N}\sum_{i=1}^{n}\left(\frac{\partial^2\log f(y_i|\theta_0)}{\partial\theta\partial\theta'}\right)\\
        & = -\hat{H}(\hat{\theta})
    \end{align*}
    \end{proof}
    \item \Emph{Outer Product of the Gradient (OPG)}: Because of the information matrix identity,
    we can also use $n^{-1}\sum_{i=1}^{n}s_i(\hat{\theta})s_i(\hat{\theta})'$.
    \begin{proof}
    \begin{align*}
        Var[s_i(\theta_0)] & = E[s_i(\theta_0)s_i(\theta)']\\
        & = \frac{1}{N}\sum_{i=1}^{N}s_i(\theta_0)s_i(\theta_0)'\\
        & = \frac{1}{N}\sum_{i=1}^{N}s_i(\hat{\theta})s_i(\hat{\theta})'
    \end{align*}
    \end{proof}
    \Info{The OPG is notorious for its poor performance with finite samples.}
\end{itemize}

\section{Regressors}


% Definition
% Theorem
% Lemma
% Example
% Proof
%\newcommand{\KeyIdea}[1]{\hltext[keybg]{\textbf{#1}}}            % 关键想法：浅绿底
%\newcommand{\Term}[1]{\textcolor{termblue}{\textbf{#1}}}           % 专业名词：品牌深绿字
%\newcommand{\Emph}[1]{\textcolor{accentD}{\textbf{#1}}}          % 强调词汇：红
%\newcommand{\Confuse}[1]{\hltext[error!20]{\textcolor{error}{#1}}} % 不懂/存疑处
%\newcommand{\Info}[1]{\textcolor{info!90!black}{#1}}             % 说明或注解